%% Lab Report for PlantillaIEEE.tex
%% V1.1
%% 2020/07/12
%% This is the template for a Lab report following an IEEE paper. Modified by Julián Valbuena after Michael Sheel original document.
%% This is a skeleton file demonstrating the use of IEEEtran.cls
%% (requires IEEEtran.cls version 1.8b or later) with an IEEE
%% journal paper.
%% Support sites:
%% http://www.michaelshell.org/tex/ieeetran/
%% http://www.ctan.org/pkg/ieeetran
%% and
%% http://www.ieee.org/
%%***************************************************************
%% Legal Notice:
%% This code is offered as-is without any warranty either expressed or
%% implied; without even the implied warranty of MERCHANTABILITY or
%% FITNESS FOR A PARTICULAR PURPOSE! 
%% User assumes all risk.
%% In no event shall the IEEE or any contributor to this code be liable for
%% any damages or losses, including, but not limited to, incidental,
%% consequential, or any other damages, resulting from the use or misuse
%% of any information contained here.
%%
%% All comments are the opinions of their respective authors and are not
%% necessarily endorsed by the IEEE.
%%
%% This work is distributed under the LaTeX Project Public License (LPPL)
%% ( http://www.latex-project.org/ ) version 1.3, and may be freely used,
%% distributed and modified. A copy of the LPPL, version 1.3, is included
%% in the base LaTeX documentation of all distributions of LaTeX released
%% 2003/12/01 or later.
%% Retain all contribution notices and credits.
%% ** Modified files should be clearly indicated as such, including  **
%% ** renaming them and changing author support contact information. **
%%*************************************************************************

\documentclass[journal]{IEEEtran}
\input{preamble}
\begin{document}
\maketitle

%%***************************************************************************
\begin{abstract}
%En esta parte, usted debe realizar una breve descripción de la esencia del documento al abordar los siguientes ítems sin mencionarlos:
%Objetivo: Exponer brevemente el problema o la cuestión abordada, en un lenguaje accesible para una audiencia científica general. Tecnología o método: Resuma brevemente la innovación tecnológica o el método utilizado para abordar el problema. Resultados: proporcione un breve resumen de los resultados y hallazgos. Conclusiones: brinde breves comentarios finales sobre sus resultados.Para una mejor comprensión de la construcción del resumen puede consultar el ejemplo que se muestra a continuación. 


\end{abstract}

\renewcommand{\IEEEkeywordsname}{Palabras Clave}
\begin{IEEEkeywords}
%% Las palabras clave, son los conceptos más relevantes tratados en su documento, procure que las palabras que aquí usted mencione sean precisas.
\end{IEEEkeywords}

\section{Introducción}
%% La introducción, a diferencia del resumen, entrega detalles respecto al contenido del documento. Es una presentación formal del informe incluyendo en lo posible referencias que soporten las afirmaciones aquí realizadas. Para una mejor comprensión de su construcción puede consultar el ejemplo que se muestra a continuación.

\IEEEPARstart{E}{n} física, 

%%*********************Desarrollo*****************************************

%%***********************Cómo Citar**************************************

\begin{comment}

Para citar en Latex, utilice el siguiente comando 
\cite{Tipler1991}

Al citar una ecuación, utilice: % (Eq.~\ref{equation:force})

Sin embargo, es importante "crear" la ecuación para poderla utilizar.
Por ejemplo: %Equation~\ref{equation:force}. 

\end{comment}

%%*******************Construir Ecuaciones*****************************

\begin{comment}

\begin{equation}
F = ma
\label{equation:force}
\end{equation}
Donde $F$ es la fuerza, $m$ es la masa y $a$ la aceleración. 

\end{comment}



%%*******************Cómo citar y construir tablas y figuras*************

%% Para citar una tabla o una figura use el siguiente comando: %~\ref{fig:ecg} or Table~\ref{table:Exps}), 

\begin{figure}[H]%[!ht]
\begin {center}
%\includegraphics[scale=0.6]{images/Overleaf.png}
%\includegraphics[width=3cm, height=3cm]{images/Overleaf.png}
%\includegraphics[scale=0.5, angle=45]{images/Overleaf.png}
\includegraphics[scale=0.5]{images/Overleaf.png}
\caption{Ejemplo de descripción de figura.}
\label{fig:ecg}
\end {center}
\end{figure}

\begin{table}[H]
\centering
\label{table:Posición vs Tiempo}
\begin{tabular}{|l|l|}
\hline
\rowcolor[HTML]{FFFE65} 
Y{[}m{]} & t{[}s{]} \\ \hline
         &          \\ \hline
         &          \\ \hline
         &          \\ \hline
         &          \\ \hline
\end{tabular} 
\caption{Ejemplo de descripción de cuadro.}
\end{table}
%you can use a table generator from here: https://www.tablesgenerator.com/#

%$^{\circ}$C


%%************************** Para usar letras griegas*********************

%Use your word processor for Greek symbols for common engineering quantities as $\beta$, $\pi$, $\gamma$ ,$\Omega$ . 



%%*************** Marco Teórico**************************************
\section{Marco Teórico}
\subsection*{A.}
{


   
}

\subsection*{B.}
{


}
\subsubsection*{1}
    { 
    
    
    }
   


%% *********************Procedimientos y Materiales***********************


\section{Procedimientos y Materiales}

%% *********************Resultados y Análisis****************************

\section{Resultados y Análisis}

%% ********************************Conclusiones*************************


\section{Conclusiones}




% **********************Referencias**********************************

% Utilice los ejemplos de referencias en el archivo example_bib.bib


\printbibliography


%********************Final del documento***********************

\end{document}


